% !PRINTED 296
\weberpage{296}

If, on the other hand, $\beta \equiv 0$ and $\alpha-\delta$ is not divisible by $p$, then
one may choose $\beta'$ in (9) arbitrarily and then determine $\alpha'$ from a congruence
 of the first degree.

If $\beta$ is not divisible by $p$, then (9) may be put into the form
\weq{10}{\left[\alpha'\beta+\beta'\left(\frac{\delta-\alpha}{2}\right)\right]^2 \equiv \beta+\beta'^2\left[\left(\frac{\alpha+\delta}{2}\right)^2-1\right] \quad (\mathrm{mod}\ p),}
and if now
\[
\text{II.} \qquad \left(\frac{\alpha+\delta}{2}\right)^2-1 \ \text{is a quadratic residue of}\ p\text{, then determine}
\]
$\beta'$ from the congruence:
\[
\beta'^2\left[\left(\frac{\alpha+\delta}{2}\right)^2-1\right] \equiv \left(\frac{\beta-1}{2}\right)^2,
\]
whereby (10) passes into
\[
\alpha'\beta+\beta'\left(\frac{\delta-\alpha}{2}\right) \equiv \pm \frac{\beta+1}{2} \quad (\mathrm{mod}\ p),
\]
and from this $\alpha'$ can be determined, whichever sign be
chosen.

If
\[
\text{III.} \qquad \left(\frac{\alpha+\delta}{2}\right)^2-1 \ \text{is a quadratic non-residue of}\ p \ \text{and at the}
\]
same time $\beta$ is a quadratic non-residue, then a non-residue
 $\nu$ can always be determined such that $\beta+\nu$ becomes a quadratic residue; for if
one lets $\nu$ run through the sequence of non-residues in $\beta+\nu$, then
the results cannot be \gesperrt{nothing but} non-residues, since among them
$\beta$ too would have to occur.

Then $\beta'$ can be determined so that
\[
\beta'^2\left[\left(\frac{\alpha+\delta}{2}\right)^2-1\right] \equiv \nu \quad (\mathrm{mod}\ p),
\]
and then $\alpha'$ from
\[
\alpha'\beta+\beta'\left(\frac{\delta-\alpha}{2}\right) \equiv \pm\sqrt{\beta+\nu} \quad (\mathrm{mod}\ p).
\]
From this it follows:

\gesperrt{In a normal divisor $\mathfrak{R}$ of the group $\mathfrak{L}_0$ a substitution
$U$ of the form}
\weq{11}{\begin{pmatrix}0, & 1\\ -1, & \xi\end{pmatrix}}
\gesperrt{must certainly occur.}

\gesperrt{Let $\xi$ first be different from zero $(\mathrm{mod}\ p)$.}
